\section{Methodology}
\label{sec:methodology}

We propose a systematic evaluation framework to assess whether generative AI fundamentally favors attackers or defenders in spam detection. Our framework consists of five principal components: data preparation through stratified sampling, synthetic spam generation with varied LLM prompt strategies, systematic mixing of real and synthetic training data, classification and comprehensive evaluation with rigorous statistical analysis. This design enables assessment of synthetic data utility across different operational scenarios while ensuring statistical rigor and reproducible conclusions.

\subsection{Data Sampling Design}

We partition the spam corpus into $R$ independent groups using stratified random sampling, with each group containing $N$ samples at fixed spam-to-ham ratio $\rho$. Sampling without replacement ensures statistical independence between groups, satisfying assumptions for parametric hypothesis testing. This replicated design enables calculation of performance variance across groups, construction of confidence intervals, and rigorous hypothesis testing with quantified uncertainty rather than point estimates.

To eliminate test set variance as a confounding factor, we maintain a fixed held-out test set across all experimental conditions. The test set remains constant while training sets vary according to synthetic proportions, ensuring that observed performance differences reflect genuine treatment effects rather than evaluation artifacts from test set composition variability. Specific values for $R$, $N$, and $\rho$ are determined through power analysis and specified in Section~\ref{sec:experiment}.

\subsection{Synthetic Data Generation}

We employ Large Language Models to generate synthetic spam through prompt-guided text generation. Given a real spam email as seed, the LLM produces a synthetic variant maintaining semantic similarity while introducing lexical and structural variation. We design three systematically varied prompt strategies to explore how generation instructions influence synthetic data quality:
\begin{itemize}
    \item \textbf{Original:} Faithfully reproduces spam characteristics while introducing natural lexical variation, maintaining semantic similarity and malicious intent without amplification or attenuation.

    \item \textbf{Strong:} Explicitly amplifies spam indicators by emphasizing urgency signals, persuasive elements, and social engineering techniques.

    \item \textbf{Weak:} Produces subtle spam with reduced explicit indicators, generating content that appears more legitimate while retaining underlying spam nature.
    
\end{itemize}

This controlled variation enables systematic investigation of whether prompt intensity affects generation quality and classifier performance. Complete prompt templates are specified in Section~\ref{sec:experiment} and Appendix~\ref{sec:appendix}.

\subsection{Cross-Group Mixing Strategy}

We employ cross-group mixing where training data from group $i$ combines with synthetic data generated from different groups $j \neq i$. This design deliberately introduces distributional mismatch between generation sources and training targets, reflecting realistic deployment scenarios where exact distribution matching between synthetic generation sources and operational training data is infeasible or unverifiable.

Cross-group mixing provides two critical advantages over within-group alternatives. First, it tests whether synthetic data provides value through semantic diversity rather than precise distribution matching, establishing whether LLM-generated spam captures generalizable features beyond specific distributional artifacts. Second, it better reflects operational constraints including temporal flexibility (using historical synthetic data for current threats), source diversity (pooling synthetic data from multiple corpora), and organizational collaboration (sharing synthetic samples without revealing operational distributions). Our preliminary analysis demonstrated that cross-group mixing maintains comparable performance to within-group mixing while providing more conservative estimates of synthetic data utility under realistic deployment conditions\footnote{Cross-group strategy shows consistent advantages over within-group strategy: mean F1 improvement of +0.009 ($p < 10^{-6}$, Cohen's $d = 0.13$) across all configurations.}.

\subsection{Evaluation Framework: Three Experimental Tracks}
Our evaluation framework encompasses three complementary tracks that systematically address the research questions outlined in Section~\ref{sec:introduction}. 


\textbf{Track 1: Synthetic-to-Real Detection (RQ1)}: This track evaluates whether LLM-generated synthetic spam can effectively train classifiers to detect real spam. Classifiers train on mixtures of real ham (fixed) and varying proportions of real spam combined with synthetic spam, then test on identical real-world data. As synthetic proportion $\theta$ varies from 0 (pure real) to 1 (pure synthetic), we quantify performance degradation patterns, establishing trade-offs between reduced real data requirements and maintained detection accuracy. The test set remains purely real throughout, ensuring results reflect operational deployment scenarios where detectors must identify genuine spam.

\textbf{Track 2a: Real-to-Synthetic Detection (RQ2a)}: This track evaluates whether classifiers trained solely on real spam can detect LLM-generated synthetic spam, representing zero-shot detection scenarios. Training employs only real ham and real spam without synthetic exposure, while testing evaluates detection of LLM-generated synthetic spam. Strong baseline performance would indicate that features distinguishing real spam from ham generalize to identifying AI-generated spam, suggesting inherent robustness. Weak baseline performance would indicate that AI-generated spam represents a novel threat requiring adaptive training strategies.

\textbf{Track 2b: Mixed-to-Synthetic Detection (RQ2b)}: This track investigates whether incorporating synthetic data from other LLMs improves detection capabilities through cross-model augmentation. Classifiers train on real ham plus varying mixtures of real spam and synthetic spam from LLM B, then test on detecting synthetic spam from LLM A (different from training). This cross-model design evaluates whether features learned from one LLM's synthetic patterns generalize to detecting another LLM's spam. Performance improvement over Track 2a baseline would demonstrate that cross-model synthetic augmentation enhances detection capabilities, enabling ``fighting fire with fire" defensive strategies.
